\section{Termination for elementary sequences with a nef b-line summand}\label{sec:threshold}

We prove termination for the elementary sequences of
Section~\ref{sec:relative} when the adjoint has an effective rational
summand and a nef rational b-line summand. Special termination
allows us to delete the coefficient-one boundary from a tail of any
purported infinite sequence.  Repeating this operation will give strictly
increasing generalized log canonical thresholds, contrary to ACC.
All programs and line-bundle comparisons take place on the full compact
spaces.  The algebraic floor-deletion argument of
\cite[Lemma~2.20 and Theorem~4.1]{CT2023} provides a useful comparison;
the relative lifts needed here will follow from the analytic results already
proved.

\begin{theorem}\label{thm:expression-termination}
Let $(X,B)$ be a rational klt pair on a globally strongly $\Q$-factorial
compact K\"ahler fourfold.  Suppose there are an effective rational divisor
$N$ on $X$ and rational nef b-line data $\bM$, carried by a projective
modification of $X$, such that
\begin{equation}\label{eq:threshold-expression}
 K_X+B\simq N+M_X
\end{equation}
as actual rational line bundles.  Then every birational elementary
$(K_X+B)$-sequence in the category of Section~\ref{sec:relative}
is finite.  In particular, this holds for every sequence constructed by
Proposition~\ref{prop:ordinary-step}, independently of its ray choices.
\end{theorem}

The nefness in this statement is nefness of the higher rational line bundle.
No global meromorphic frame for that bundle, or uniform Cartier index for
its traces on the successive spaces, is assumed.

\subsection{Thresholds and their lifts}

The first lemma identifies the threshold on a fixed resolution.  The
negative curve ensures that the testing boundary on that resolution is
nonzero, which makes the threshold finite.

\begin{lemma}\label{lem:rational-threshold}
Let $V$ be a globally strongly $\Q$-factorial compact K\"ahler space,
let $B,N\geq0$ be rational divisors, and let $\bM$ be rational nef b-line
data on a projective higher model.  Set
\[
 T=N+M_V,\qquad A(t)=K_V+B+tN+tM_V.
\]
Suppose $(V,B+aN+a\bM)$ is generalized klt for some rational $a\geq0$,
and $T\cdot C<0$ for a compact curve $C\subset V$.  Then
\[
 s=\sup\{t\geq0:(V,B+tN+t\bM)\text{ is generalized lc}\}
\]
is a finite rational number strictly greater than $a$.  The pair is
generalized lc at $s$ and generalized klt at every $0\leq t<s$.
\end{lemma}

\begin{proof}
Choose a projective log resolution $p:W\to V$ carrying $\bM$ and resolving
all the divisor comparisons below.  Write
\[
 K_W+C_W=p^*(K_V+B),\qquad
 H=p^*(N+M_V)-M_W.
\]
These are the comparisons of actual rational sheaves from
Lemma~\ref{lem:trace}.  Strong $\Q$-factoriality makes $N$ and $M_V$
rational Cartier, and
\begin{equation}\label{eq:threshold-defect}
 H=p^*N+(p^*M_V-M_W)\geq0
\end{equation}
by Lemma~\ref{lem:negativity}.  After increasing $W$, the support of
$C_W+H$ is simple normal crossing.  The crepant higher boundary at
parameter $t$ is $C_W+tH$.

The divisor $H$ is nonzero.  Otherwise $p^*T=M_W$ is nef.  Since $p$ is
projective, there is a compact curve $\widetilde C\subset W$ mapping onto
$C$ with positive degree: take a component dominating $C$ in $p^{-1}(C)$
and intersect it with sufficiently many general relative hyperplanes.
The projection formula would give
\[
 0\leq M_W\cdot\widetilde C
   =\deg(\widetilde C/C)\,T\cdot C<0.
\]
Here only degrees on curves are used; no absolute curve criterion for
analytic nefness is needed.

For the finitely many prime components $E$ with
$h_E=\operatorname{coeff}_E H>0$, put
$c_E=\operatorname{coeff}_E C_W$.  The log-smooth discrepancy criterion
gives
\begin{equation}\label{eq:rational-threshold-formula}
 s=\min_{h_E>0}\frac{1-c_E}{h_E}.
\end{equation}
Indeed all coefficients of $C_W+aH$ are strictly below one, and components
with $h_E=0$ impose no new bound.  Formula~\eqref{eq:rational-threshold-formula}
proves finiteness, rationality, attainment, and $s>a$, as well as the
claimed singularities at and below $s$.
\end{proof}

At the threshold we need a generalized dlt model to use special
termination.  The next lemma lifts each step of an infinite sequence to
a finite, nonempty string on such models.  It uses the existence of a
finite relative generalized klt run from
Proposition~\ref{prop:relative-finite}; all comparisons concern one flip
base at a time.

\begin{lemma}\label{lem:threshold-lift}
Consider an infinite sequence of small elementary transformations
\[
 X_i\dashrightarrow X_{i+1},\qquad
 X_i\xrightarrow{f_i}Z_i\xleftarrow{f_i^+}X_{i+1},
\]
between globally strongly $\Q$-factorial compact K\"ahler spaces of
dimension at most four.  Let
$B_i,N_i\geq0$ be strict transforms of rational divisors, let $\bM$ be
fixed rational nef b-line data, and put
\[
 T_i=N_i+M_{X_i},\qquad A_i(t)=K_{X_i}+B_i+tT_i.
\]
Suppose that for rational numbers $0\leq a<s$:
\begin{enumerate}[label=\textup{(\roman*)}]
 \item $(X_i,B_i+aN_i+a\bM)$ is generalized klt and
       $(X_i,B_i+sN_i+s\bM)$ is generalized lc for every $i$;
 \item $T_i$ and $A_i(t)$ have negative degree on the contracted ray
       and their transforms are relatively ample on the positive side,
       for every rational $t\geq a$.
\end{enumerate}
Then the sequence lifts to a concatenation of nonempty finite strings of
elementary generalized dlt steps at parameter $s$.  At the ends of the
strings there are projective crepant morphisms
\[
 h_i:(Y_i,\Delta_i+s\bM)\longrightarrow
       (X_i,B_i+sN_i+s\bM)
\]
with effective generalized dlt boundary, extracting only divisors of
coefficient one.
\end{lemma}

\begin{proof}
Take $h_1$ by Proposition~\ref{prop:gdlt-mod}.  Suppose the current $h_i$
has been constructed.  We choose a nearby klt parameter, run a finite
relative program, and identify its endpoint with a crepant model of
$X_{i+1}$.

\smallskip\noindent
\emph{Choosing a nearby klt parameter.}
For $t\in[a,s]$, define its crepant boundary by
\[
 K_{Y_i}+\Delta_i(t)+tM_{Y_i}=h_i^*A_i(t).
\]
This boundary depends affinely on $t$.  Its nonexceptional coefficients
are $b+tn$ with $b,n\geq0$.  Its exceptional coefficients equal one at
$s$.  There are only finitely many exceptional primes, so for a rational
$t_i<s$ sufficiently close to $s$, with $t_i\geq a$, the boundary
$\Delta_i(t_i)$ is effective.  In particular, a nonexceptional coefficient
zero at the positive parameter $s$ has $b=n=0$ and stays zero.  The pair
downstairs is generalized klt for $a\leq t<s$: its discrepancies are
affine combinations of the positive discrepancies at $a$ and the
nonnegative ones at $s$.  Thus the crepant pair at $t_i$ upstairs is
effective and generalized klt.

Choose a curve in the ray of $f_i$.  By the one-ray condition on global
rational line-bundle degrees, there is a positive rational number $q_i$
such that $A_i(t_i)-q_iA_i(s)$ has zero degree on all $f_i$-contracted
curves.  Apply Lemma~\ref{lem:descent}, using the generalized klt
antiample adjoint at $t_i$, to obtain a rational line bundle $L_i$ on
$Z_i$ with
\begin{equation}\label{eq:threshold-proportionality}
 A_i(t_i)\simq q_i A_i(s)+f_i^*L_i.
\end{equation}
After pullback by $h_i$, this is an actual identity modulo a rational
bundle from $Z_i$, not merely an equality of numerical classes.

\smallskip\noindent
\emph{Running the finite relative program.}
The composite $Y_i\to Z_i$ is projective and bimeromorphic.  Apply
Proposition~\ref{prop:relative-finite} to the effective generalized klt
pair at $t_i$ to obtain some finite relative elementary run ending nef
over $Z_i$.  Transforming~\eqref{eq:threshold-proportionality} through
each step shows that the same run is negative, and positive after each
flip, for the full adjoint at $s$.  It is therefore a generalized dlt
run at $s$, by Proposition~\ref{prop:relative-steps}, and its endpoint
is nef for this adjoint as well.

\smallskip\noindent
\emph{Identifying the endpoint.}
To check the endpoint, on a common resolution write
\begin{equation}\label{eq:threshold-comparison}
 P_0=P_j+G_j=P_++F.
\end{equation}
Here $P_0$ is the pullback of $h_i^*A_i(s)$, $P_j$ is the pullback of the
lifted endpoint adjoint, and $P_+$ is the pullback of $A_{i+1}(s)$.
Negativity for the lifted run gives $G_j\geq0$, exceptional over its
endpoint, and negativity for the original small step gives $F\geq0$,
exceptional over $X_{i+1}$.  Both endpoint adjoints are nef over $Z_i$,
and $A_{i+1}(s)$ is ample there.  Lemma~\ref{lem:nef-end-comparison}
applies.  Explicitly, if $p$ and $q$ denote the maps to these two targets,
then $P_j-P_+=F-G_j$ is $q$-nef with pushforward $-q_*G_j\leq0$;
its negative is $p$-nef with pushforward $-p_*F\leq0$.  The two
negativity inequalities give $P_j=P_+$.  Relative ampleness and
projective connected fibers then give the crepant projective morphism
from the lifted endpoint to $X_{i+1}$.

Every prime exceptional for this morphism is the transform of an
$h_i$-exceptional prime: the original step is small and the lifted run
extracts none.  If it survives, its coefficient one is preserved by the
boundary pushforward rule.  This proves the required assertion for
$h_{i+1}$.  The string is nonempty, since $h_i^*A_i(s)$ has negative
degree on a curve lifting an $f_i$-contracted curve and so is not nef
over $Z_i$.  Induction proves the lemma.  The choice of $t_i$ is made
separately at each finite stage.
\end{proof}

\subsection{Deleting the floor and raising the threshold}

\begin{proof}[Proof of Theorem~\ref{thm:expression-termination}]
From an infinite sequence we will construct successively larger
thresholds with fixed coefficient data.  At each repetition,
special termination will let us remove the floor while retaining an
infinite sequence with the required signs.

\smallskip\noindent
\emph{Step 1: the infinite tail and its first threshold.}
Suppose there is an infinite sequence.  By Lemma~\ref{lem:cycle-loss},
it has only finitely many divisorial steps.  Truncate after the last of
them and relabel the resulting sequence by $X_1\dashrightarrow X_2
\dashrightarrow\cdots$.  Taking actual reflexive traces transports
\eqref{eq:threshold-expression} through this finite prefix.  The
effective divisor remains effective and rational, and a projective
common higher model carries the pullback of the original nef line
bundle.  Thus the relabelled input still has the required expression.

We will repeat the following construction.  Its input is an infinite
small elementary sequence with rational families
\begin{equation}\label{eq:threshold-family}
 (X_i,B_i+tN_i+t\bM),\qquad
 A_i(t)=K_{X_i}+B_i+tN_i+tM_{X_i},
 \quad T_i=N_i+M_{X_i},
\end{equation}
and a rational $a\geq0$ such that the pairs at $a$ are generalized klt,
and both $T_i$ and $A_i(t)$ have the negative-to-positive signs of each
step for all rational $t\geq a$.  The divisors $B_i,N_i$ are effective
strict transforms.  Initially $a=0$, the nef part at zero is zero, and
\[
 T_i\simq K_{X_i}+B_i,
 \qquad A_i(t)\simq(1+t)(K_{X_i}+B_i),
\]
so these conditions hold.

Let $s$ be the threshold of the first pair in~\eqref{eq:threshold-family}.
Lemma~\ref{lem:rational-threshold} applies to a contracted curve and gives
$s\in\Q$, $a<s<\infty$.  Since $A_i(s)$ is negative on the input side
and ample on the positive side at each step, discrepancy monotonicity
shows that every pair in the sequence at $s$ is generalized lc.
Lemma~\ref{lem:threshold-lift} gives an infinite concatenation of
generalized dlt elementary steps at $s$, with compatible crepant
endpoint morphisms to the $X_i$.

\smallskip\noindent
\emph{Step 2: deleting the floor after special termination.}
Apply Lemma~\ref{lem:cycle-loss} and Theorem~\ref{thm:special} to this
concatenation.  It has a tail consisting of small steps whose exceptional
loci, on both sides, avoid the support of the coefficient-one boundary.
Every lifted string is finite and nonempty, so an endpoint occurs beyond
any prescribed finite prefix.  Start the tail at such an endpoint
$h:Y\to X_i$ and relabel.  Write the union of the supports of $B_i,N_i$
as $\bigcup P$, with coefficients $b_P,n_P\geq0$.  If $\widehat P$
denotes the strict transform on $Y$, the crepant boundary at $s$ is
\[
 \Delta(s)=\sum_P(b_P+sn_P)\widehat P+\sum_{E\text{ exceptional for }h}E.
\]
Set
\begin{equation}\label{eq:retained-threshold-coefficients}
 \widetilde B=\sum_{b_P+sn_P<1}b_P\widehat P,
 \qquad
 \widetilde N=\sum_{b_P+sn_P<1}n_P\widehat P.
\end{equation}
Then
\[
 \widetilde B+s\widetilde N
    =\Delta(s)-\lfloor\Delta(s)\rfloor.
\]
Deleting the floor of a generalized dlt pair gives a generalized klt
pair with the same nef data.  Indeed every place of discrepancy zero
has center in that floor, whose rational Cartier pullback has positive
order at the place; all other discrepancies were already positive.
Consequently $(Y,\widetilde B+s\widetilde N+s\bM)$ is generalized klt.
Both divisors in~\eqref{eq:retained-threshold-coefficients} are effective.

\smallskip\noindent
\emph{Step 3: retaining the signs and increasing the threshold.}
Let their strict transforms define the new family on every model of
the lifted tail.  Its signs require checking for the testing direction,
as well as for its adjoint.  Do this over each original flip base $Z_i$.
At the initial endpoint of its lifted string, the trace comparison gives
\begin{equation}\label{eq:floor-supported-test}
 h_i^*T_i=\widetilde N_i+M_{Y_i}+F_i,
 \qquad
 \Supp F_i\subseteq\Supp\lfloor\Delta_i(s)\rfloor.
\end{equation}
To see the support assertion, the difference between pullback and strict
transform of $N_i$, and the difference between pullback and b-line trace
of $M_{X_i}$, are $h_i$-exceptional.  Every such prime has coefficient
one in $\Delta_i(s)$.  The remaining difference consists of the discarded
$n_P\widehat P$, also supported in the floor.  This is a comparison of
actual rational sheaves; a sign for $F_i$ is unnecessary.

On $X_i$, both $A_i(s)$ and $T_i$ have negative degree on the contracted
ray, and their ratio is a positive rational number $c_i$.
Lemma~\ref{lem:descent}, applied using any nearby generalized
klt parameter below $s$, gives a rational line bundle $Q_i$ on $Z_i$
such that
\begin{equation}\label{eq:floor-test-proportionality}
 h_i^*A_i(s)\simq
 c_i(\widetilde N_i+M_{Y_i}+F_i)+(f_i h_i)^*Q_i.
\end{equation}
Transform this identity through the finite lifted string.  The tail is
small, so all its boundary divisors, traces, and $F_i$ transform strictly.
The support of the transform of $F_i$ remains in the transformed floor.
All contracted and flipped curves are disjoint from that floor.
Equation~\eqref{eq:floor-test-proportionality} therefore shows that
\(\widetilde T=\widetilde N+M_{\mathrm{trace}}\) has negative degree on
every contracted ray and positive relatively ample transform on the
other side.  For the ampleness assertion, the same identity, with the
floor bundles trivial near the exceptional fibers, identifies its
restriction there with a positive rational multiple of the full ample
adjoint; fiberwise ampleness is the relative ampleness criterion.

Likewise subtracting the floor changes neither sign of the adjoint at
$s$.  On each side, for every rational $u\geq s$, the new adjoint is
\[
 \widetilde A(u)=\widetilde A(s)+(u-s)\widetilde T,
\]
so it has the same negative-to-positive signs.  At later endpoints the
retained coefficients still obey~\eqref{eq:retained-threshold-coefficients}:
the downstairs maps are small, and an exceptional divisor of an endpoint
morphism cannot become the transform of a downstairs prime.  All
boundaries upstairs are pushforwards, hence strict transforms on this
tail.  Thus the families constructed over consecutive bases agree.

We have obtained another infinite small elementary sequence satisfying
the same conditions, now with $a=s$.  The generalized klt condition at
$s$ holds throughout the tail by discrepancy monotonicity.  It also
implies the generalized klt condition at parameter zero, by
\eqref{eq:threshold-defect}.  Lemma~\ref{lem:rational-threshold} shows
that its next threshold is finite, rational, and strictly larger than
$s$.  The construction can therefore be repeated indefinitely if the
original sequence was infinite.  It produces
\begin{equation}\label{eq:increasing-thresholds}
 0<s_1<s_2<s_3<\cdots.
\end{equation}

\smallskip\noindent
\emph{Step 4: the fixed data for analytic ACC.}
Let $I$ be the finite set consisting of zero and the coefficients of the
first $B$, and let $J$ consist of zero, the coefficients of the first
$N$, and $1/p$, where $p>0$ makes $pM_{W_0}$ integral on the
original nef carrier $W_0$.  At every repetition, the new boundary and testing
coefficients are obtained solely by omissions from the previous ones,
as in~\eqref{eq:retained-threshold-coefficients}.  Newly extracted
coefficient-one primes are removed before the next threshold input.
The boundary coefficients therefore stay in $I$, the testing-boundary
coefficients stay in $J$, and the nef part at parameter zero remains
zero.  On every required higher model the testing nef data are
$\frac1p$ times the pullback of the same integral nef line bundle.
The denominator $p$ is fixed.  Common projective models are needed only
for finitely many maps at a time and do not change $p$.

We verify explicitly that each $s_j$ belongs to the analytic generalized
threshold set with these data.  Choose a general point of a computing
center on its compact input, and a relatively compact Stein neighborhood
$U$ of that point.  Restrict a projective log resolution carrying the
data to $U$, and denote it by $g:W\to U$.  The computing place still
meets this restriction.  Hence the local threshold equals $s_j$: the
pair is lc at $s_j$, whereas the same place has negative discrepancy
for every larger parameter.

Let $\cL$ represent the integral line bundle $pM_W$.  The coherent sheaf
$g_*\cL$ has rank one on the isomorphism locus of $g$.  Cartan's
Theorem~A supplies a section nonzero at a point of that locus, and thus
a nonzero section of $\cL$ on $W$.  Its divisor $P'$ is Cartier and
represents $\cL$.  In particular $P'$ is $g$-nef and $M_W$ is represented
by $\frac1pP'$.  The induced meromorphic identification gives precisely
the original trace and pullback defect.  Compatible local canonical
representatives may be chosen in the same way from the coherent
canonical sheaf.  Neither choice changes the boundary or testing-boundary
coefficients, the discrepancies, or the factor $1/p$.  The testing
boundary is rational Cartier because its components are restrictions
of global rational Cartier divisors on the strongly $\Q$-factorial input.

The analytic generalized threshold theorem
\cite[Theorem~1.1 and Theorem~3.3; see also Definition~2.3 and
Remark~2.5]{HXACC2023} consequently applies in dimension four with the
fixed DCC sets $I,J$.  Its nef summands are actual relatively nef Cartier
divisors with the fixed weight $1/p$; it requires no common bound for
Cartier indices of their downstairs traces.  It prohibits the strictly
increasing sequence~\eqref{eq:increasing-thresholds}.  This contradiction
proves the theorem.
\end{proof}
